\title{X-LoRA: Mixture of Low-Rank Adapter Experts, a Versatile Framework for Large Language Models with Applications in Protein Mechanics and Molecular Design}

\begin{abstract}
We present a mixture of expert approach to develop fine-tuned large language models using a deep, layer-wise token-level method grounded in low-rank adaptation (LoRA). Beginning with a collection of pre-trained LoRA adapters, our gating mechanism leverages hidden states to dynamically blend adapted layers, enabling the X-LoRA model to utilize diverse capabilities and generate novel deep layer-wise combinations to address various tasks. This design draws inspiration from biological principles of universality and diversity, wherein neural network components are repurposed in different hierarchical configurations. As a result, the X-LoRA model can be seamlessly integrated with any existing large language model (LLM) without requiring structural modifications. We create a specialized X-LoRA model endowed with scientific functionalities, including forward and inverse analysis tasks alongside improved reasoning, with a focus on biomaterial analysis, protein mechanics, and design. The significance of this work lies in providing access to easily scalable and adaptable models that possess robust domain expertise and the ability to synthesize knowledge across disciplines. Incorporating experts in biology, mathematics, reasoning, bio-inspired materials, mechanics and materials science, chemistry, protein biophysics, mechanics, and quantum mechanics-based molecular properties, we conduct a range of physics-centric case studies. These investigations cover knowledge retrieval, protein mechanics forward/inverse problems, protein design, adversarial agentic modeling including ontological knowledge graph construction, and molecular design. The model not only delivers quantitative predictions for nanomechanical properties of proteins and quantum mechanical molecular characteristics but also performs reasoning over these outcomes, accurately predicting plausible mechanisms that explain distinct molecular behaviors.
\end{abstract}

\section{Introduction}  
Large language models (LLMs)~\cite{Vaswani2017AttentionNeedc,Touvron2023LlamaModels,OpenAI2023GPT-4Report,Chowdhery2022PaLM:Pathways,Jiang2023Mistral7B,Gunasekar2023TextbooksNeed} have become widely popular, including in the creation of specialized models that excel in particular tasks, reasoning processes, or scientific disciplines~\cite{Bubeck2023SparksGPT-4,Buehler2023MechGPTModalities,Nejjar2023LLMsAnalysis,Buehler2023GenerativeDesign,Luu2023BioinspiredLLM:Materials,Luu2023GenerativeSolvents,Buehler2023MeLMProblemsb,Ge2023OpenAGI:Experts}. Nevertheless, training these models can be expensive, especially when a broad range of capabilities is desired. Approaches such as low-rank adapters (LoRA)~\cite{hu2021lora} have been introduced as more resource-efficient alternatives, though these adaptations typically concentrate on more specialized knowledge areas. The core idea behind LoRA is to introduce low-rank matrices added to the original full-rank weight matrices, designating these low-rank matrices as the sole trainable parameters. Because only the adapter layers are updated during training, these models generally retain the knowledge embedded in the pre-trained base model while tailoring the model to particular tasks, all with lower computational cost. The efficiency of training LoRA adapters enables the development of numerous specialized models using this method.

Although LoRA adapters offer an effective approach to building specialized model functions, a key open problem is how to combine multiple such adapters into enhanced models that provide a unified and potentially superior set of capabilities. In fact, integrating LLMs has emerged as an active research direction~\cite{Kim2023SOLARUp-Scaling}. Experimental investigations, including approaches such as spherically interpolated model parameters (SLERP) and related techniques~\cite{Arcee-ai/mergekit:Models.}, have demonstrated promising outcomes, indicating the potential to merge models and unlock new functionalities. Models combined using these strategies are generally computed {\it a priori} following specific algorithms and, while they have shown encouraging results on some benchmarks, further study is needed to systematically design mixed models and understand how particular features arise.

We propose that dynamically mixing parameters represents a more versatile technique, as it enables the combined model to explore novel parameter combinations during inference.

Similar methods have employed mixture of experts (MoE) frameworks~\cite{Jacobs1991AdaptiveExperts,Hampshire1992TheRecognition,Jordan1994HierarchicalAlgorithm}, including the recently introduced Mixtral LLM~\cite{Jiang2024MixtralExperts}. Nevertheless, these approaches often entail high computational costs and substantial resource requirements even during inference. In this work, we introduce a computationally efficient mixture-of-expert method that utilizes a collection of distinct LoRA adapters, which can be trained independently with ease. Our method employs a fine-grained and adaptable mechanism, inspired by a biological framework that reuses neural network components to build a hierarchy of interacting models, ranging from foundational models to agentic systems (see Fig.~\ref{fig:hierarchical_concept}). To demonstrate its quantitative effectiveness, we concentrate on applying this strategy specifically to science-oriented LLMs, with a particular focus on protein mechanics within the broader domain of bio-inspired materials analysis and design.

We present an algorithm that dynamically mixes adapters at the token-level state, allowing for fine-grained mixing of adapters across individual layers. This approach, named X-LoRA, provides access to a diverse range of functionalities essential for scientific applications of multimodal language modeling~\cite{Hu2023DeepScience,Buehler2023AMetamaterials,Buehler2022ProteinNetworks}.

\subsection{Fundamental concepts of X-LoRA}

We begin with a brief overview of the principles underpinning LoRA adapters. The core idea behind low-rank adaptation (LoRA)~\cite{hu2021lora} posits that weight updates have a low “intrinsic dimension,” leveraging this by freezing the original weights $W_0 \in \mathbb{R}^{d \times k}$ and restricting updates to a low-rank decomposition

\begin{equation}
W_0 + \Delta W_0 = W_0 + BA 
\end{equation}

where $B \in \mathbb{R}^{d \times r}$ and $A \in \mathbb{R}^{r \times d}$, with the rank $r \ll {\rm min}(d,k)$. During training with LoRA~\cite{hu2021lora}, the pretrained weights $W_0$ remain fixed, and only the matrices $A$ and $B$ are updated as trainable parameters. The forward propagation through a LoRA layer can be described by the following equation:

\begin{equation}
    h = W_0x + \Delta W_0x = W_0x + BAx
\end{equation}

It is important to note that, in practice, LoRA adapters include a fixed scalar multiplier $\alpha$. This scalar is applied to the decomposed matrices, resulting in

\begin{equation}
    h = W_0x + BAx \cdot \alpha
\end{equation}

Building on this concept, X-LoRA introduces the idea of employing a collection of LoRA adapters, each possessing distinct capabilities, as demonstrated in prior scientific LLM research (e.g., \cite{Buehler2023GenerativeDesign}). Instead of statically combining or integrating models, we propose a dynamic gating mechanism that scales each individual LoRA adapter with fine-grained control over tokens and layers to enable mixing deep within the model. More specifically, for a given LoRA adapter $i$, the original parameter $\alpha_i$ is adjusted by a scaling factor denoted as $\lambda_i$, producing a new scaling factor $\alpha^*_i$ for the adapter:

\begin{equation}
    \alpha^*_i = \alpha_i \cdot \lambda_i
\end{equation}

This scaling factor $\lambda_i$ is predicted by an X-LoRA scaling head, which leverages the model’s hidden states and constitutes the trainable element within the X-LoRA framework. Because the scaling head can utilize complex encoded representations, it can be efficiently realized as a simple feed-forward neural network composed of one or a few layers (users of X-LoRA can specify the architectural details of this neural network).

The scaling process is performed for every token in the input sequence. Consequently, this provides a high degree of granularity, allowing each LoRA adapter active at particular layers within the LLM to be individually scaled.

From this point onward, we refer to this mechanism of controlling the adapters’ components in this manner as scaling. Additional details regarding this method are outlined in the Materials and Methods section.

\subsection{Paper outline}

The structure of this paper is outlined as follows. Initially, we describe the methodology and training approach, including the relevant datasets that contribute to the creation of an X-LoRA model possessing specialized capabilities in the physical sciences, with a particular emphasis on biomaterials such as proteins. Subsequently, we showcase a series of experiments where the X-LoRA model is applied to multiple tasks, including question answering, conversational and agentic modeling, protein design and analysis, among others. We perform an in-depth examination of the scaling behaviors and validate our approach by comparing it with molecular modeling and other physical data and methodologies.

\section{Results and discussion}

The development of the X-LoRA model follows these steps:

\begin{enumerate}
    \item Train a foundational base LLM (completed in prior research) \item Independently train a collection of adapters to gain specialized knowledge in various domains (these represent distinct or overlapping skill and knowledge areas) \item Train the integrated X-LoRA model
\end{enumerate}

Our experimental phase begins with the training of several LoRA adapters. We create nine adapters, each fine-tuned with unique expertise, based on the Zephyr-7B-$\beta$ model~\cite{HuggingFaceH4/zephyr-7b-betaFace}, which itself is built upon the Mistral-7B model\cite{Jiang2023Mistral7B}:

\begin{enumerate}
    \item Bioinspired materials \item Chain-of-thought (CoT) and reasoning \item Chemistry \item Mathematics \item Physics \item Biology \item Mechanics and materials \item Platypus: Logic and reasoning \item Protein mechanics tasks (generative sequence-to-property and inverse capabilities)
\end{enumerate}

These particular fields of concentration and expertise were selected because our ultimate goal is to develop a STEM-oriented, materials science-centered LLM. Nonetheless, we hypothesize that our methodology can be broadly applied to a wide range of domains, particularly to generate modalities that facilitate the formation of connections between diverse knowledge areas~\cite{Buehler2023MechGPTModalities,Buehler2023MeLMProblemsb}.

Following the creation of the LoRA adapters, we proceed to train the X-LoRA model. For this purpose, we utilize a subset derived from the combined collection of training samples originally employed for training the individual LoRA adapters (further specifics are provided in the Materials and Methods section, including detailed descriptions of the dataset, implementation, and related information).

We treat the X-LoRA system as an experimental tool to evaluate its effectiveness across various scenarios. We carry out a series of tests encompassing complex question answering, application of the generative model to protein design, analysis and reasoning tasks, as well as developing more advanced interaction techniques through adversarial agentic modeling. Throughout these evaluations, we compare the performance of the base model against that of the X-LoRA model, while monitoring how the scaling head engages different layers of the various experts during the generation or formulation of responses.

\subsection{Question answering and observed X-LoRA layer-wise scaling weights}

Initially, we investigate the distinctions between the X-LoRA model and the base foundational model. We pose the following question to both models:

\begin{quote}
\small\texttt{What material behavior (stiffening vs softening) causes supersonic crack speeds in mode II cracks? Think step by step.}
\end{quote}

The base model answers:

\begin{quote}
\small\texttt{Supersonic crack speeds, also known as Rayleigh waves, occur in mode II cracks due to the stiffening material behavior of the material near the crack tip. In mode II fracture, the loading is shear-dominated and the crack propagates perpendicular to the direction of applied stress. Here's a step-by-step explanation:

As the crack tip nears the material, the stresses around it escalate rapidly. This creates an area of intense stress concentration, known as the stress singularity.

Within this stress singularity, the material undergoes elastic deformation that causes an increase in stiffness, referred to as stiffening behavior. This stiffening occurs because the material close to the crack tip becomes increasingly constrained as the crack advances [...]}
\end{quote}

Drawing on our own expertise~\cite{Buehler2008AtomisticFailure}, this explanation is not only inaccurate but also unnecessarily verbose. By contrast, the X-LoRA model provides the following response:

\begin{quote}
\small\texttt{The material behavior that causes supersonic crack speeds in mode II cracks is stiffening.}
\end{quote}

This response is accurate, concise, and straightforward~\cite{Buehler2008AtomisticFailure}.

The distinction between the base model and the X-LoRA model is even more pronounced for specialized protein analysis and design tasks. The task presented is:

\begin{quote}
\small\texttt{CalculateForceHistory< M S K S N T Y R M L V L L E D D T K I N K E D E K F L K G K P G K M H E F V D E L I L P F N V D E L D E L N T W F D K F D A E I C I P Y V E I L K E S G M K > }
\end{quote}

The response from the base foundation model is:

\begin{quote}
\small\texttt{I'm sorry but the input provided is not a valid mathematical expression. It appears to be a sequence of characters that do not form a calculable function. Please provide a valid mathematical expression for me to calculate.

However, I can tell you that the sequence of characters you provided is actually a series of function and variable names that may be part of a larger program or script in a specific programming language. Without context, it's impossible to determine the intended calculation.

If you provide more information, such as the programming language and the context in which this expression appears, I may be able to provide more insight into its meaning and possible calculations. }
\end{quote}

In contrast, the X-LoRA model replies as follows, successfully accomplishing the design task of producing the force-deformation curve linked to this protein:

\begin{quote}
\small\texttt{[0.013,0.066,0.078,0.045,0.072,0.126,0.127,0.107,0.119, 
0.074,0.106,0.098,0.103,0.128,
0.077, 0.069,0.131,0.108,0.144,0.174, 0.166,0.183,0.185,0.217,0.175,0.230,0.187,0.224, 0.267,0.301,0.315]}
\end{quote}

Figure~\ref{fig:QA_weights_sample_0} presents outcomes from this question answering alongside the observed X-LoRA scaling weights. The two queries, one addressing mechanics and materials (Fig.~\ref{fig:QA_weights_sample_0}a) and the other protein mechanics (Fig.~\ref{fig:QA_weights_sample_0}b), exhibit distinct activation patterns. Notably, we observe that a complex arrangement of scaling values is employed in each scenario, implying that the X-LoRA model exploits a heterogeneous mixture of different adapters across layers. Additionally, the heatmap in Fig.~\ref{fig:QA_weights_sample_0}a displays a wider range of activation spanning many experts and layers, with certain regions showing strong activation. Conversely, the heatmap in Fig.~\ref{fig:QA_weights_sample_0}b reveals a more focused activation pattern, with fewer regions of high intensity. We interpret this as suggesting that the gating mechanism on the right acts more selectively in engaging experts.

In Fig.~\ref{fig:QA_weights_sample_0}, the bright line-like regions indicate paths of high activation, demonstrating that the scaling function selectively activates certain experts more strongly at specific layers. The observed patterns suggest that importance is not evenly spread across all experts and layers; rather, it is typically sparse and selective. This sparsity aligns with findings in other mixture-of-experts models, where only a limited subset of experts is engaged for each input, specializing in distinct areas of the problem domain.

By examining the maps across different tasks, both the sparsity and the specific activation patterns may be crucial for gaining insight into which experts the model depends on for various input types or during different stages of processing within the neural network (we will later discuss how these maps evolve dynamically during the resolution of complex tasks).

To deepen this analysis, we now investigate a broader range of questions to understand how the X-LoRA model employs its various experts when presented with diverse prompts. We are particularly interested in whether, and in what manner, the X-LoRA model combines different adapters when questions involve overlapping skills or domains. Fig.~\ref{fig:QA_weights} presents results from question answering alongside observed X-LoRA scaling weights, with the questions given in Table~\ref{tab:table_qa}. Table~\ref{tab:table_qa} provides responses from both X-LoRA and the base model, Zephyr-7B-$\beta$. The contrasts between the two models are notable.

To clearly illustrate this, incorrect or questionable answers are highlighted in \redhighlight{red}. The comparison demonstrates that X-LoRA delivers significantly higher accuracy and more concise responses.

As anticipated, we detect complex adapter mixing and frequent activation of multiple dominant LoRA experts. For example, Fig.~\ref{fig:QA_weights}d shows that a question regarding pine cone mechanics prompts strong activation of the mechanics/materials adapter along with the CoT/reasoning and bioinspired adapters. Although X-LoRA selects the correct letter choice, the answer is partially incorrect since pine cones open in dry, not wet conditions. However, a follow-up question clarifying the release conditions yields the correct answer.

In a protein-related question, Fig.~\ref{fig:QA_weights}f reveals adapter mixing in response to an inquiry about the mechanics of Bombyx mori silk, activating a combination of the bioinspired, mechanics/materials, and biology adapters. For highly specialized tasks such as protein design (Fig.~\ref{fig:QA_weights}a), we observe primarily singular activation of a specific expert.

A more detailed examination of the heatmap of scalings shown in Fig.~\ref{fig:QA_weights}a reveals that, aside from the clearly prominent importance of the protein mechanics expert, there are multiple bands where certain experts receive considerable weight across several layers. This pattern is especially evident for the mechanics/materials and reasoning experts. These bands are discontinuous and exhibit a fluctuating pattern, which could indicate that these experts are particularly effective under specific conditions or for particular features. Notably, the reasoning expert (positioned immediately to the left of the protein mechanics expert) is distinguished by having the highest activation across several layers, particularly in the upper layers (approximately layers 100 to 140). We suggest this may point to a crucial influence in the model’s final output or decision process, though further investigation is required to confirm this. Additionally, there is a localized but strong activation of the physics and biology experts at certain layers, potentially reflecting a specialized role for these experts in processing features or representations within those regions.

Moreover, Fig.~\ref{fig:QA_token_history} presents a token-level history reflecting the cumulative scalings across all layers. This analysis highlights the histories for two example cases: in Fig.~\ref{fig:QA_token_history}a, a task concerning pine cone seed release, and in Fig.~\ref{fig:QA_token_history}b, the outcome of a sequence of protein mechanics tasks. The protein mechanics task involved a series of prompt-answer interactions, starting with two property calculations and culminating in a generative task to design a new protein. Although expert usage remains generally consistent, notable shifts occur in the engagement of the LoRA experts during the process. In particular, the protein mechanics expert is predominantly activated in the final stages of generation, when the new protein is being created.

Table~\ref{tab:table_qa} presents various use cases spanning multiple domains along with a comparison to the performance of the base model alone. We explore several of these cases in greater detail.

For example, we ask a question that combines elements of protein science, biology, and logic, focusing on protein expression patterns. X-LoRA provides the correct answer. In contrast, the base model's response contains certain inaccuracies and logical errors when evaluating the marker protein expressions in each cell type. The base model accurately notes that fibroblasts express Protein B and do not express Protein A due to the mutual exclusivity between Protein A and Protein B. However, it fails to explicitly acknowledge that fibroblasts also express Protein C. Since there is no rule preventing the co-expression of Protein C with Protein B, and fibroblasts have no limitation against Protein C, it logically follows that fibroblasts express Protein C as well. Regarding Neurons, the base model correctly begins by stating that neurons express Protein A and not Protein B, following the mutual exclusivity rule. Yet, the explanation about whether neurons can express Protein C or both Protein C and Protein B is confusing and contains an error. Since neurons express Protein A and not Protein B, the only other possibility, according to the rules, is that they express Protein C. There is no scenario in which neurons express Protein B, contrary to what the initial analysis suggests. Hence, neurons express Protein A and Protein C, but not Protein B. For Epithelial cells, the base model rightly observes that epithelial cells do not express Protein A but do express both Protein B and Protein C. Nevertheless, the rationale provided is unclear and includes incorrect claims about mutual exclusivity between Protein A and Protein C, which was not part of the original assumptions. The straightforward logic is that epithelial cells do not express Protein A but must express the other two proteins, Protein B and Protein C, because the initial guidelines only specify that Protein A and Protein B cannot be co-expressed.

In contrast, X-LoRA accurately determines the marker proteins expressed by each cell type based on the given observations and logical inferences. For fibroblasts, X-LoRA correctly infers that fibroblasts do not express Protein A because the presence of Protein A excludes Protein B expression. Since fibroblasts do express Protein B, and considering that each cell type expresses a unique combination of three marker proteins (an incorrect assumption introduced here, as the original statement did not specify the number of proteins each cell type expresses, only that they have a unique combination), the conclusion that fibroblasts also express Protein C is valid given the information provided (excluding the misunderstanding regarding the number of proteins). Regarding neurons, the model correctly identifies them as expressing Protein A and Protein C. The reasoning is logical: Protein B cannot co-express with Protein A, and since neurons express Protein A and one other marker protein (not Protein B), they must express Protein C. Lastly, for epithelial cells, X-LoRA correctly concludes that they express Protein B and Protein C, because epithelial cells do not express Protein A, and by elimination, given that they express two marker proteins, these must be Protein B and Protein C.

When addressing the question about pine cone opening under humid versus dry conditions, the base model provides a confused and factually incorrect response. Conversely, X-LoRA delivers a clear and well-reasoned answer. Other related examples include a question on the exceptional mechanical properties of Bombyx mori silk, where the base model incorrectly claims that silk fibers conduct electricity. Additionally, the base model fails in various domain-specific tasks, such as the question regarding the role of hyperelasticity in maximum fracture speeds, where X-LoRA offers a concise and accurate response, while the base model does not provide a valid answer. In a query about the strong adhesion of gecko feet, X-LoRA correctly highlights the significance of setae and nanoscale effects, whereas the base model wrongly attributes adhesion to special glue proteins. These and other examples, discussed in the table and throughout the paper, strongly demonstrate the superior capabilities of X-LoRA compared to the base model.

Figure~\ref{fig:perf_comp_bioinspiredexam} presents a comparison of knowledge recall evaluation experiments using the bio-inspired knowledge test set introduced in \cite{Luu2023BioinspiredLLM:Materials}. It is evident that X-LoRA achieves the highest performance, despite being a considerably smaller model than the BioinspiredLLM or Llama-BioLLM models (7B parameters for X-LoRA versus 13B parameters for the others). Additionally, several other comparisons with recently released LLMs, including Orca-13B and Llama-13b-chat, are provided.

The figure also compares X-LoRA with the base Zephyr-7B-$\beta$ model, demonstrating a substantial improvement in performance with X-LoRA. These results suggest that excellent performance can be attained using X-LoRA even when smaller base models are employed.

Figure~\ref{fig:image_synth} illustrates an application in image synthesis, where X-LoRA is used to create a prompt for other models such as DALL-E 3, as shown here.

\begin{quote}
\small\texttt{Create a prompt that I can use for image generation that accurately describes the microstructure of a hierarchical composite.

Explicitly describe the geometry.}
\end{quote}

The comparison in Figure~\ref{fig:image_synth}(b)-(c) clearly demonstrates that a more concise prompt leads to a significantly more accurate image generated by DALL-E 3~\cite{DALLECard}. Notably, the prompt produced by Zephyr-7B-$\beta$ fails to accurately capture the given instruction and introduces extraneous design elements such as fibers and nanoparticles, which were not part of the original prompt focused solely on the microstructure of a hierarchical composite. In contrast, the prompt generated by X-LoRA is shorter and more focused, resulting in a much better output.

\subsection{Generative solutions to protein analysis}

We now turn to specialized generative and analytical tasks related to proteins. These capabilities arise in the X-LoRA model through the protein mechanics adapter and are further strengthened by integrating the model’s knowledge in biology, bio-inspired materials, reasoning, and logical deduction. This section focuses exclusively on protein-related tasks, evaluating the quantitative performance of this functionality.

Figure~\ref{fig:protein_force_history_prediction} presents results derived from the protein task aimed at predicting force-deformation behaviors based on sequences. Figure~\ref{fig:protein_force_energy_prediction} illustrates the overall accuracy in forecasting unfolding force and unfolding energy from the sequence, comparing actual values with predictions (yielding $R_2=0.85$). In general, the model demonstrates excellent performance and, as will be discussed in the following section, also performs exceptionally well on the inverse task of design and cycle consistent evaluation\cite{Lu2023GenerativePropertiesb}. However, unlike previous studies that relied on highly specialized GPT-type models to tackle these tasks, our X-LoRA model integrates these specialized capabilities alongside a variety of other deep specializations introduced through multiple fine-tuned adaptations.

\subsection{Protein design and analysis}

We now broaden the scope of test cases to more integrative applications that engage multiple areas of expertise. Specifically, the experiments presented in this section illustrate how the X-LoRA model can leverage a unique combination of skills and how agentic modeling can yield even more complex and comprehensive outcomes. We start the discussion with a protein design task, where the model is tasked with designing a novel protein to achieve a specific target force-deformation response. These force-deformation responses have been explored in prior research~\cite{Ni2023ForceGen:Model} and represent measurements of force as a function of deformation when the protein is stretched at both ends (corresponding to a classic protein unfolding experiment~\cite{Lu1998UnfoldingSimulation.}).

Fig.~\ref{fig:design_protein} illustrates the application of the generative protein task. We design a protein exhibiting a specified force-deformation response (the training data and boundary conditions are derived from original molecular dynamics (MD) simulations of protein pulling, see~\cite{Ni2023ForceGen:Model}), and subsequently evaluate the predicted sequence to determine if it aligns with the intended target properties (Fig.~\ref{fig:design_protein}a). This cycle-consistent evaluation is crucial for verifying that the inverse design capabilities correspond with the forward prediction.

Fig.~\ref{fig:design_protein}b presents the sequence predicted: \texttt{MSKSNTYRMLVLLEDDTKINKEDEKFLKGKPGKMHEFVDELILPFNVDELDELNTWFDKFDAEICIPYVEILKESGMK}. Additionally, Fig.~\ref{fig:design_protein}c displays a folded structural model of the protein, generated via AlphaFold 2~\cite{Jumper2021HighlyAlphaFoldb}. To investigate the relationship of the designed protein with known sequences, Fig.~\ref{fig:design_protein}d analyzes it using a Basic Local Alignment Search Tool (BLAST) Tree~\cite{Altschul1990BasicTool}. Notably, the designed protein shows some connections to hypothetical proteins previously examined in earlier research but which have not yet been structurally characterized or identified.

Employing agentic modeling and integrating multiple X-LoRA models as components in a multi-agent framework, we conduct a further analysis of the predicted sequence with our X-LoRA model. This leverages the enhanced ability of our model to comprehend protein mechanics alongside its integration with other knowledge domains. Using the aforementioned prediction as input, we request an in-depth analysis and design recommendations for sequence variants. The task is initiated as follows:

\begin{quote}
\small\texttt{I have identified this amino acid sequence: MSKSNTYRMLVLLEDDTKINKEDEKFLKGKPGKMHEFVDEL ILPFNVDELDELNTWFDKFDAEICIPYVEILKESGMK.

How can I use it?

Please provide logical steps, specifically addressing the engineering aspects of the problem.}
\end{quote}

We employ multiple agentic interactions to develop the solution, applying adversarial strategies (details in Materials and Methods) to encourage the model to probe deeper and explore more intricate aspects of the concepts under discussion. The summary of this dialogue, as identified by the X-LoRA model (the full conversation is provided in the Supplementary Material, SI Conv Example 1), is as follows:

\begin{quote}
\small\texttt{The dialogue explores the possible uses of a particular amino acid sequence in the field of protein engineering. This sequence comprises both hydrophobic and hydrophilic residues, indicating its potential to form a stable conformation and engage with other molecules. To design a protein incorporating this sequence, computational approaches can predict its secondary and tertiary structures; however, challenges include the intricate nature of protein folding and the associated computational expense. Enhancements can be achieved by creating more efficient algorithms and integrating novel experimental findings. It is essential to validate computational models to guarantee their precision and practical applicability. Methods such as cross-validation, regularization, early stopping, ensemble techniques, and transfer learning are employed to avoid overfitting and underfitting. Additionally, understanding and explaining these models can aid in guiding real-world implementations.}
\end{quote}

We also request the model to generate a summary of main points, resulting in:

\begin{quote}
\small\texttt{\begin{itemize}
    \item The amino acid sequence MSKSNTYRMLVLLEDDTKINKEDEKFLKGKPGKMHEFVDELI LPFNVDELDELNTWFDKFDAEICIPYVEILKESGMK holds promise in protein engineering because of its combination of hydrophobic and hydrophilic residues. Computational tools can forecast the protein’s secondary and tertiary structures, though obstacles such as protein folding complexity and computational demands exist.
    \item Ensuring computational model validation is vital for accuracy and effective practical use. Techniques including cross-validation, regularization, early stopping, ensemble methods, and transfer learning are applied to mitigate overfitting and underfitting.
    \item Model interpretation and explanation also contribute to enhancing practical applications.
\end{itemize}
}
\end{quote}

Ultimately, the main conclusion is:

\begin{quote}
\small\texttt{The key takeaway is that the amino acid sequence MSKSNTYRMLVLLEDDTKINKEDEKFLKGKPGKMHEFVDELILPFNVDELDELNTWFDKFDAEICIPYVEILKESGMK holds promise for protein engineering due to its combination of hydrophobic and hydrophilic residues. To design a protein with this sequence, computational approaches can be employed to predict its secondary and tertiary structures; however, challenges include the complexity of protein folding and the computational resources required. Ensuring the accuracy and practical utility of computational models necessitates validation, using techniques such as cross-validation, regularization, early stopping, ensemble approaches, and transfer learning to avoid overfitting or underfitting. Additionally, interpreting and explaining the models can aid in guiding practical implementations.}
\end{quote}

We can further delve into detailed sequence engineering strategies. For this, we ask the model to provide more targeted design recommendations aimed at enhancing the protein’s stability:

\begin{quote}
\small\texttt{How could I modify the following protein sequence to improve its stability: MSKSNTYRMLVLLEDDTKINKEDEKFLKGKPGKMHEFVDELILPFNVDELDELNTWFDKFDAEICIPYVEILKESGMK? }
\end{quote}

A summary of the discussion (see Supplementary Material, SI Conv Example 2 for the full transcript) is:

\begin{quote}
\small\texttt{The dialogue centers on strategies to increase protein stability by incorporating hydrophobic amino acids and forming disulfide bonds. The biologist explains that hydrophobic residues typically contribute to protein stability, and disulfide bonds can further enhance this stability. The engineer requests concrete examples of hydrophobic amino acids and guidance on the number of disulfide bonds to introduce. The biologist cites examples of proteins successfully stabilized via these modifications, such as insulin, bovine pancreatic trypsin inhibitor, and green fluorescent protein. The conversation also addresses possible drawbacks of adding hydrophobic residues or disulfide bonds, including potential alterations in protein function, added complexity, and risks of aggregation. The biologist notes these issues can be mitigated through careful amino acid selection, computational evaluation, and experimental validation.}
\end{quote}

The summary of the main points:

\begin{quote}
\small\texttt{\begin{itemize}
\item To increase the stability of a protein sequence, one can incorporate hydrophobic amino acids and disulfide bonds.

Hydrophobic amino acids generally contribute to enhancing protein stability, and disulfide bonds also improve structural stability.

Notable proteins that have been successfully stabilized by these modifications include insulin, bovine pancreatic trypsin inhibitor, and green fluorescent protein.
\item Possible drawbacks of adding hydrophobic amino acids or disulfide bonds include alterations in protein function, added complexity, and a tendency toward aggregation.
\item These challenges can be mitigated by careful amino acid selection, computational modeling, and experimental confirmation.
\end{itemize}
}
\end{quote}

The essential takeaway is:

\begin{quote}
\small\texttt{The most important conclusion from the discussion is that introducing hydrophobic amino acids and disulfide bonds can enhance the stability of the protein sequence MSKSNTYRMLVLLEDDTKINKEDEKFLKGKPGKMHEFVDELILPFNVDELDELNTWFDKFDAEICIPYVEILKESGMK. This directly addresses the initial query.}
\end{quote}

From the standpoint of protein engineering, these recommendations are logical as they would lead to designs with improved stability. Specific examples and approaches are outlined that researchers may further investigate.

In another case, we analyze three sequences. Initially, we request the model to evaluate the stability of each, followed by asking it to identify the strongest candidate and justify why that protein is likely the most stable. Additionally, we ask for an explanation of the laboratory synthesis steps. The three sequences (manually derived by modifying one sequence from the test set) are:

\begin{quote}
\small\texttt{
\begin{itemize}
    \item LNTWFDKFDAEICIPYVEILKESGMK \item AITWFDKFDAEICIPYVEALKIAAMK \item AAAAAAAAKFDAAEICIIAAAAAAAAKIAAMK
\end{itemize}
}
\end{quote}

This scenario evaluates whether the model can integrate various capabilities and perform this adaptively throughout the conversation. In other words, to fluidly switch between different skills such as protein mechanics, logical reasoning, hypothesis generation, and so forth.

The dialogue is summarized in Figure~\ref{fig:conversation_evolve}, alongside an illustration of how the scalings evolve throughout the conversation (the analysis presents both the deep layer-wise adapter scaling and a combined assessment to provide a metric for the effective utilization of the different experts). The findings clearly indicate that the protein task is initially highly pertinent, progressively shifting towards a stronger representation of the bioinspired adapter and the biology adapter. By the conclusion of the conversation, the biology adapter emerges as the most dominant.

To evaluate the predictions and assess the quality of the responses, the three proteins examined in this task are shown in Figure~\ref{fig:protein_visuals} (proteins folded using Alpha Fold 2~\cite{Jumper2021HighlyAlphaFoldb}, through ColabFold \cite{Mirdita2022ColabFold:All}). We observe that the most stable structure in the center displays the most well-organized secondary structure, consistent with the idea that it produces the highest unfolding force, as predicted by the X-LoRA model in Figure~\ref{fig:conversation_evolve}. The other two structures are less organized; the first one is primarily helical with some turns, while the last one is partially disordered and features a kink in its geometry. These characteristics account for the lowest unfolding force predicted by the model.

By examining the visual depictions of the folded configurations, we verify that the central structure, which is the most stable, exhibits the most well-organized secondary arrangement. This aligns with the concept that it produces the greatest unfolding force as forecasted by the model. We observe that the other two models are less structured: the first is predominantly helical with some turns, while the last is partially disordered and contains a kink in its geometry. We propose that these characteristics likely account for the lowest observed unfolding force predicted by the model, consistent with the predictions of the X-LoRA model.

In a second example, we concentrate on the energetic attributes of proteins, specifically the unfolding energy as the protein is stretched from its initial conformation to a fully unfolded state~\cite{Ni2023ForceGen:Model}. We again analyze three sequences (the first two are the same as those previously considered, but the third is a longer construct that combines the second sequence at the start and end with the first sequence in the middle):

\begin{quote}
\small\texttt{\begin{itemize}
    \item LNTWFDKFDAEICIPYVEILKESGMK \item AITWFDKFDAEICIPYVEALKIAAMK \item AITWFDKFDAEICIPYVEALKIAAMKLNTWFDKFDAEICIPYVEILKESGMKAITWFDKFDAEICIPYVEALKIAAMK
\end{itemize}
}
\end{quote}

Initially, we request the model to determine the unfolding energy for each sequence, and then to interpret the results: based on the amino acid sequences, we ask it to explain why one protein is probably the most stable. Finally, we inquire about the potential functions that the most stable protein might possess.

Figure~\ref{fig:MJB_20} presents a transcript of a dialogue that engages multiple specialized experts and combines knowledge from these domains, demonstrated here by analyzing the unfolding energy of three proteins. By the conclusion of the discussion, the CoT/reasoning adapter becomes the dominant component as the X-LoRA model addresses a question requiring reasoning over the outcomes to predict likely protein structural characteristics as well as possible protein functions. In particular, the mechanistic inquiry regarding the basis of protein stability highlights the formation of hydrophobic interactions among nonpolar amino acids along with hydrogen bonding between polar amino acids. The model forecasts the involvement of several hydrophobic amino acids (I, F, W, Y) and hydrogen bonding amino acids (D, E, K) in enhancing protein stability.

Additionally, the model anticipates that the presence of cysteine (C) residues in the sequence would facilitate the creation of disulfide bonds, which further reinforce the protein’s structure. Notably, the critical predictions concerning structural elements explaining the high stability are accurately identified, as confirmed by structural analyses depicted in Fig.~\ref{fig:MJB_21}. This figure’s analysis corroborates the existence of all these features. It also clarifies why the other two proteins, consisting of much shorter segments with variable helical domain stability, exhibit significantly lower unfolding energy.

\subsection{Adversarial agentic modeling to connect distinct scholarly disciplines and knowledge yielding ontological knowledge graph generation}

Given that our X-LoRA model possesses extensive cross-domain expertise, we utilize it to explore links among disparate concepts, knowledge repositories, and fields of specialization. In one experiment, we pose two queries to the model, each crafted similarly but with a different emphasis. Subsequently, we generate triplets to construct an ontological knowledge graph that organizes the responses into more structured outputs~\cite{Giesa2012CategoryDesign,Spivak2011ReoccurringAnalogies}. The resulting graph delivers a cohesive understanding of the derived insights and visually represents the relationships between concepts in an interpretable and mechanistic fashion.

We employ two adversarial agents to address the question. Agent 1 is labeled as a "critical physicist," while Agent 2 takes the role of a "philosopher." Agent 1 initiates the inquiry and is responsible for probing and generating further questions in response to each answer provided by Agent 2. As the dialogue progresses, the discussion delves more deeply into the topic and reveals various facets. The question is deliberately designed to encompass multiple disciplinary domains to test the model’s capacity to synthesize ideas and concepts.

The question posed is:

\begin{quote}
\small\texttt{Consider an abstract representation of a classical musical composition as a graph of interacting notes, chords, melodies, rhythms and so on.

What happens if we would apply pressure, up to a point of failure, on this abstract concept of music?

Provide logical steps, mathematical reasoning, and discuss specifically the physics of the problem.}
\end{quote}

The entire conversation is available in the Supplementary Material (refer to SI Conv Example 3). For brevity, we present here only the summary and the key insights, capturing the main points:

\begin{quote}
\small\texttt{The dialogue explores the effect of applying pressure to an abstract representation of a classical musical composition, conceptualized as a physical system consisting of interacting notes, chords, melodies, and rhythms. Pressure can be quantified through stress, with the critical failure point identified when the stress reaches a threshold value. To validate the mathematical model’s predictions, a physical analog of the abstract musical composition can be employed, and experimental results compared with the theoretical forecasts. Limitations of the proposed model include the simplified depiction of interactions, static nature, assumption of homogeneous distribution, and linear stress-pressure relationship. Enhancing the model’s fidelity requires incorporating more detailed and realistic interactions, accounting for dynamic effects, introducing heterogeneity, and including non-linear terms.

[...]

The principal takeaway is that applying pressure to an abstract representation of a classical musical composition can be interpreted as altering the force field between the interacting musical elements. As pressure escalates, so does the system’s internal stress, ultimately reaching a critical point where the system collapses under its own influence. This understanding provides a physical rationale for how pressure impacts the abstract notion of music and how this effect can be measured using stress.}
\end{quote}

Although the generated text offers valuable insights, structuring the data into an ontological representation can facilitate more interpretable analysis~\cite{Buehler2023MechGPTModalities}. For graph construction, we merge the entire conversation and generate a knowledge graph. An example graph is depicted in Figure~\ref{fig:graph}. The resulting graph is partitioned into communities via the Girvan-Newman algorithm~\cite{Girvan2002CommunityNetworks}, identifying a total of 12 communities. The figure displays both the overall graph (Figure~\ref{fig:graph}a) and detailed views with node labels (Figure~\ref{fig:graph}b-c).

\subsection{Creation of X-LoRA-Gemma Incorporating Combined Protein, Chemical, Bio-inspired, and Mechanics of Materials Functionalities}
\label{gemma-section}

To demonstrate that our proposed methodology is applicable to other foundational models with different architectures, we trained an alternative X-LoRA model based on the Gemma-7B-it model~\cite{Google/gemma-7b-itFace}. For this model, we implemented a set of four LoRA adapters defined as follows:

\begin{enumerate}
    \item Bioinspired materials \item Mechanics and materials \item Protein mechanics tasks (which include generative sequence-to-property and inverse modeling capabilities) \item Quantum mechanics-based molecular properties QM9 (featuring generative SMILES-to-property and inverse design functionalities; refer to Table~\ref{tab:table_QM9} for definitions and a summary of all properties calculated within that dataset)~\cite{Ruddigkeit2012EnumerationGDB-17,Ramakrishnan2015ElectronicSpace}
\end{enumerate}

The X-LoRA-Gemma model was developed similarly to the original X-LoRA model but is designed to encompass distinct functional capabilities, particularly enhancing the ability to predict 12 quantum mechanical properties as well as to design molecules that satisfy these 12 quantum mechanical property criteria. This new adapter was trained alongside the others listed above, which utilize the same datasets employed in our previous experiments. Note that while most of this paper centers on the Zephyr-based X-LoRA model, for clarity in this section we differentiate between the models by referring to the original as X-LoRA-Zephyr and the Gemma-based variant as X-LoRA-Gemma.

Figure~\ref{fig:protein_force_energy_prediction} illustrates the overall performance of the X-LoRA model, showing predictions of unfolding force and unfolding energy from sequence data ($R_2 = 0.87$), as well as prediction accuracy for quantum mechanical molecular properties derived from the QM9 dataset ($R_2 = 0.96$). We observe outstanding results for both tasks, with the protein design task even performing slightly better than the original X-LoRA model.

Next, we showcase the capabilities of X-LoRA-Gemma in a straightforward multi-agent design scenario. We examine a small organic molecule, \texttt{CCCC12CC1C(=O)N2} (extracted from the QM9 test set). This bicyclic compound contains a ketone (=O) and an amine (N) integrated into the ring; see Figure~\ref{fig:QM9_design}(a) for a depiction of the molecule. Initially, we prompt the X-LoRA-Gemma model with:

\begin{quote}
\small\texttt{Here is a molecule in SMILES code: CCCC12CC1C(=O)N2

What modifications would I need to make to render this molecule more suitable for polymer design?}
\end{quote}

The response is:

\begin{quote}
\small\texttt{To enhance the utility of this molecule for polymer development, you might:
\begin{enumerate}
\item  Add a reactive functional group, like an amine or hydroxyl group, to enable further modification and crosslinking.
\item  Extend the carbon backbone to be longer and more flexible, improving the molecule's capability to form polymers with diverse structures and characteristics.
\item  Attach a side chain containing a functional group, such as a carboxyl group, to facilitate covalent bonding with other molecules and polymers.
\item  Integrate a cyclic structure, for instance a lactam ring, to produce a polymer that is more rigid and stable.
\item  Incorporate a variety of side chains to generate a wide range of properties and functionalities within the polymer.
\item  Vary the number of carbons in the backbone to enable the synthesis of polymers with different molecular weights and features.
\item  Introduce a crosslinking site within the molecule to form a polymer that is more stable and less reactive.
\item  Embed a repeating unit structure to design a polymer with specific molecular weight and characteristics.
\item  Combine several of the above strategies to develop a versatile and functional polymer suitable for multiple applications.
\end{enumerate}}
\end{quote}

Next, we concentrate on the first recommendation.

\begin{quote}
\small\texttt{How can I enhance the overall reactivity of a molecule? I intend to use this molecule for constructing a polymer material.

The molecule under consideration, represented by SMILES, is: CCCC12CC1C(=O)N2

I have the option to adjust the following properties:

Dipole moment, $\mu$ Polarizability, $\alpha$ Highest Occupied Molecular Orbital energy, HOMO Lowest Unoccupied Molecular Orbital energy, LUMO HOMO-LUMO gap: The energy difference between HOMO and LUMO, gap Electronic spatial extent Zero-point vibrational energy Heat capacity at constant volume Internal energy at 0 K Internal energy at 298.15 K Enthalpy at 298.15 K Free energy at 298.15 K Which are the top THREE properties I should modify, and in what direction (increase or decrease) should these changes be made?}
\end{quote}

The reply from X-LoRA-Gemma is:

\begin{quote}
\small\texttt{\begin{itemize}
    \item  HOMO-LUMO gap (decrease) \item  Dipole moment (increase) \item  Polarizability (increase)
\end{itemize}
Altering these properties can enhance the overall reactivity of the molecule, making it better suited for use in polymer material synthesis. }
\end{quote}

We employ both the generative and predictive functions of X-LoRA-Gemma to achieve these targeted changes in molecular characteristics.

For the optimization, we implement a two-agent framework consisting of two X-LoRA-Gemma instances: one agent carries out the generative task \texttt{GenerateMolecularProperties<...>} to propose candidate molecules. The second agent then performs a calculation using \texttt{CalculateMolecularProperties<...>} to evaluate the difference between the desired and predicted molecular properties. This cycle continues until either the maximum iteration count is reached or the mean-squared error falls below a predefined threshold.

Figure~\ref{fig:QM9_design} presents the outcomes of this procedure. Panel (a) illustrates the targeted design alongside the deviation from the original molecule's properties. Panel (b) displays the top 20 molecules generated by the algorithm, with panel (c) showing the distribution of mean-squared error, and panel (d) plotting the mean-squared error across the generated designs. The optimal design identified is \texttt{CC(C)C1=CC(=O)NC=C1}, which closely aligns with the target criteria as depicted in panel (e). Lastly, panel (f) provides an optimized 3D molecular structure, where we applied UFF~\cite{Rappe1992UFFSimulations} for structural optimization.

To evaluate the molecule’s reactivity further, we calculated the Gasteiger charges~\cite{Gasteiger1980IterativeCharges} (shown in Figure~\ref{fig:QM9_design}(f) for visualization). The molecule contains an aromatic ring featuring a ketone (C=O) group, which is more electrophilic than typical carbon-carbon or carbon-hydrogen bonds, thus prone to nucleophilic attack. The combination of the nitrogen-containing aromatic ring and the ketone forms a pyridone ring. This structure exhibits increased reactivity compared to a simple aromatic ring due to the heteroatom (nitrogen) and the carbonyl group, creating reactive sites for chemical processes such as nucleophilic addition at the carbonyl carbon or electrophilic substitution on the ring. The nitrogen atom, as part of the heteroaromatic system, influences the molecule’s aromaticity and reactivity by modifying electron density distribution across the ring. This adjustment affects the molecule’s interactions with other chemicals, making the nitrogen-containing ring a focal site for reactions, particularly electrophilic attacks facilitated by the electron-donating character of the nitrogen atom.

The carbon atom within the ring, positioned adjacent to both an oxygen and a nitrogen atom and bearing a partial positive charge (+0.25), implies that it is somewhat electron-deficient. This occurs because oxygen, being more electronegative, draws electron density toward itself, resulting in the carbon having reduced electron density. This effect is partially offset by the neighboring nitrogen, which is less electronegative than oxygen but more so than carbon. Nonetheless, the existence of a partial positive charge indicates that this carbon is a probable site for nucleophilic attack, where a nucleophile (an electron-rich species) would be attracted to this electron-poor carbon.

The nitrogen atom, exhibiting a partial negative charge (-0.33), suggests it is comparatively electron-rich relative to its typical state. This can be attributed to its lone pair of electrons and the influence of the aromatic ring's electron system. In aromatic heterocycles, the nitrogen's lone pair may participate in the delocalized $\pi$-electron system, though the precise extent of this contribution depends on the structure of the ring and the electronegativity of neighboring atoms. The partial negative charge indicates that nitrogen is more nucleophilic, implying a greater tendency to donate electrons, either by forming bonds with electrophiles or undergoing protonation. Additional quantum mechanical analyses are recommended to explore this further.

The carbon atom with a partial positive charge represents a reactive site for nucleophilic attacks. Other molecules may interact with this site via mechanisms where nucleophiles target electron-deficient carbons.

The hydrogen atom bonded to the nitrogen atom, which carries a partial positive charge of 0.17, is well-suited for forming hydrogen bonds with other atoms or molecules. Hydrogen bonds are a type of dipole-dipole interaction occurring when a hydrogen atom bonded to a highly electronegative atom (such as nitrogen, oxygen, or fluorine) interacts with another electronegative atom possessing a lone pair of electrons. This characteristic renders the molecule potentially capable of acting as a hydrogen bond donor, a crucial aspect affecting its solubility in water and other polar solvents, as well as its interactions within biological systems.

The ketone group located next to the nitrogen-containing aromatic ring imparts distinctive reactivity that can be harnessed in the synthesis of polymers~\cite{McQuarrieSimonPhysicalChemistry,CareySundbergAdvancedOrganicChemistry,Varghese2022BeyondApplications,Varghese2022BeyondApplications}. This compound is classified as a lactam, due to the cyclic amide structure formed by the ketone and nitrogen within the ring. Lactams play a significant role in polymer chemistry, especially in the creation of polyamides, a class of polymers widely utilized in various material applications. The ketone functional group is capable of undergoing several chemical reactions valuable in polymer chemistry, such as nucleophilic addition. This reaction can be employed to lengthen the polymer chain or to attach side chains that may alter the polymer’s characteristics. The nitrogen atom in the aromatic ring can markedly influence the polymer’s electronic attributes and engage in hydrogen bonding, which might affect the polymer’s melting temperature, solubility, and mechanical behavior. N-heteroaromatic compounds are also recognized for their capacity to coordinate with metal ions, enabling potential applications in fields like catalysis or the development of materials with electronic or photonic functionalities. Additionally, this molecule could participate in polymerization through its amide (lactam) group; polyamides formed in this way are noted for their durability, flexibility, and resistance to wear and chemicals, making them suitable for a broad array of uses ranging from textiles (e.g., nylon) to automotive components~\cite{Varghese2022BeyondApplications,Varghese2022BeyondApplications}.

Alkyl substituents can enhance the hydrophobic nature of the molecule, so we expect that polymers synthesized from monomers containing these alkyl groups will exhibit increased hydrophobicity, influencing their solubility across various solvents as well as their interaction with water. This property could be beneficial for uses that demand water resistance or particular solubility traits. Additionally, the incorporation of alkyl groups can influence the polymer's mechanical characteristics. For example, they may function as plasticizers within the polymer matrix, thereby improving the polymer's flexibility. Depending on the alkyl groups' length and branching, this might result in materials that are more flexible or possess superior impact resistance.

\subsubsection{Additional benchmarking} We present further benchmarking results of the X-LoRA-Gemma model, depicted in Figure~\ref{fig:perf_MMLU}, utilizing a subset of the MMLU benchmark centered on chemistry, physics, and biology (fields closely related to those on which our model was trained)~\cite{Hendrycks2020MeasuringUnderstanding}. Our findings indicate that both X-LoRA models examined in this study outperform their respective base models. Specifically, the X-LoRA models exceed the base models’ performance (Zephyr model: 54.6\% for the base versus 56.6\% for the X-LoRA variant; Gemma model: 40.6\% for the base and 52.4\% for the X-LoRA-Gemma model). The highest overall accuracy is achieved by the X-LoRA-Zephyr model, though both versions are comparatively close in performance.

\section{Conclusion}

Multimodal LLMs have numerous current and prospective applications in science, accompanied by a significant demand for specialized, domain-specific expertise. Beyond general language processing, the ability to perform computational or generative tasks, as demonstrated here with protein mechanics, is crucial~\cite{Luu2023BioinspiredLLM:Materials,Buehler2023MechGPTModalities,Buehler2023MeLMProblemsb}. X-LoRA, which is based on open-source language models, allows for the dynamic blending of expertise and knowledge from different fields, as illustrated in several examples including Figure~\ref{fig:conversation_evolve}. This study revealed how the X-LoRA model autonomously and adaptively evolves its use of adapters, mixing them in a sophisticated manner across the LoRA layers. Such complex integration of LoRA layers is consistently observed in all cases examined here, as shown in Figures~\ref{fig:QA_weights_sample_0} and \ref{fig:QA_weights}. This raises an intriguing challenge to better comprehend the fundamental mechanisms from a physics standpoint, where an LLM can be viewed as a large-scale graph-forming system that leverages complex graph structures. Consequently, the mixing of adaptation experts at deeper layer levels, as executed here, modifies these mechanisms in ways that foster the development of effective solutions.

The algorithm introduced enhances the traditional inference approach of a single forward pass by incorporating two forward passes. This trains the model to initially `reflect' on the question and consider how to reconfigure itself before producing a response. This represents a straightforward implementation of ‘self-awareness,’ enabling the model to adjust its own architecture to optimally address a task. As a result, despite having a relatively modest size of 7 billion parameters, the model can reason across multiple scientific fields (such as biological materials, mathematics, physics, chemistry, logic, mechanics, etc.), greatly improving its ability to generate novel solutions. Additionally, this approach may inspire the design of alternative model architectures, for example, methods that extend the reconfiguration strategy to parts or even the entirety of non-adapted systems. In such cases, the scaling head would dynamically reconfigure sections of a trained model or combine multiple models. As demonstrated with X-LoRA, these strategies can serve as powerful tools to broaden or integrate the capabilities of distinct models.

We conducted multiple performance evaluations, including a comparison with significantly larger models (Figure~\ref{fig:perf_comp_bioinspiredexam}), where X-LoRA demonstrated clear superiority. The comprehensive comparison presented in Table~\ref{tab:table_qa} revealed that X-LoRA consistently outperforms the base model across a range of tasks and application areas. Additional comparisons involved quantitative evaluations of numerical prediction tasks, such as protein mechanics and quantum-mechanical molecular properties, where X-LoRA achieved high accuracy with $R_2$ values ranging from 0.85 to 0.96. For context, these results exceeded other reported outcomes, such as protein unfolding force prediction using the MeLM model~\cite{Buehler2023MeLMProblemsb}, which reported an $R_2$ of 0.78. The ability to accurately handle multiple prediction tasks, alongside inverse design challenges and advanced reasoning skills, was highlighted in the molecular design example (Section~\ref{gemma-section}). This example further demonstrated that X-LoRA can be effectively applied using different base model architectures (Gemma versus Zephyr/Mistral). Future research could extend this work by developing X-LoRA models for even larger base models.

A limitation of the proposed method is the requirement for two forward passes, which can increase computational overhead: one to compute hidden states used as inputs to the X-LoRA scaling head, and another with $\Lambda$ applied to the adapters. However, the first pass benefits the X-LoRA scaling head by leveraging inherent LLM capabilities, in contrast to separately learning the knowledge through training a larger scaling head. For practical implementation in model-serving systems, distinct key-value caches need to be maintained for both scaling and forward passes. To enhance inference speed, we created \texttt{Mistral.rs}, a LLM serving platform implemented in Rust~\cite{matsakis2014rust} that supports X-LoRA and incorporates multiple performance optimization techniques. Despite this, a key advantage of our approach is its applicability to any model without necessitating changes to internal logic, which we consider beneficial, particularly given its compatibility with the extensive resources in the Hugging Face ecosystem.

A key benefit of the proposed method is its straightforward implementation in any existing LLM (e.g., with the release of the X-LoRA code, it should be readily applicable to any model within the Hugging Face ecosystem). Another benefit lies in how the scaling head leverages the inherent strengths of the employed LLM, learning intricate ways to optimally combine adapters for particular responses. When training X-LoRA, it is advantageous to include complex question-answer or conversational samples that enable the model to discover the most effective combinatorial approach. We conducted several analyses related to protein design and mechanics. For example, Figure~\ref{fig:MJB_20} and Fig.~\ref{fig:MJB_21} illustrate stability assessments of multiple protein sequences, featuring comparisons, reasoning over the predicted outcomes, and mechanistic explanations that provide a fundamental chemical basis for the predicted behavior. As shown in Fig.~\ref{fig:MJB_21}, this was corroborated through a structural examination of the protein’s folded 3D conformation.

Regarding future directions, further investigation could focus on expanding the diversity of adapter experts beyond the initial expertise represented in our LoRA set. The X-LoRA model outperforms both the base model and the base model equipped with a single adapter, delivering responses to complex queries by drawing on specialized knowledge embedded in the adapter set. There also appear to be synergistic interactions among different adapters, as suggested by the intricate mixing patterns of scaling weights across layers. This phenomenon merits additional study and could be compared with approaches like SLERP, which integrate models through alternative strategies. While we trained the X-LoRA scaling head using a subset of the training data—comprising a few hundred samples from each original dataset that informed the adapter development—a more targeted training dataset tailored for this purpose could be designed. For instance, we observed that multiple-choice questions tend to activate reasoning-focused adapters, likely because they were trained on such data. Additional prompting is necessary to help the scaling head grasp specific domains relevant to posed questions, possibly through particular system messages or other techniques. Overall, the creation of appropriate training sets is a promising area for further exploration, including methods to effectively induce layer-wise scaling mixing. Our experiments tracking these mechanisms over extended conversations demonstrate that the model can dynamically adjust scaling to optimally address different tasks, which is an encouraging outcome.

We identified intriguing patterns regarding how experts are activated throughout the layers, with noteworthy observations elaborated in the paper. This analysis can be extended and might provide valuable understanding about the benefits and mechanisms of mixing model components or layer sets. Such exploration would be especially compelling when applying the X-LoRA method beyond protein mechanics and design, as demonstrated here. We reserve this exploration for future research.

Another significant direction is the deeper examination of agentic modeling, exemplified here by the interaction of two adversarial X-LoRA models. Subsequent studies could involve more than two models to enrich interactions and introduce additional functionalities, thereby steering models into novel generative domains. This approach could also facilitate the creation of techniques that integrate physics-based modeling or other validation procedures, employing code-writing/executing agents or agents utilizing function calling or other processing methods~\cite{Ghafarollahi2024ProtAgents:Learning}.

\section{Materials and Methods}

The X-LoRA codes and illustrative examples can be found at \url{https://github.com/EricLBuehler/xlora}. Model weights and datasets related to the X-LoRA discussed here, along with further examples, are accessible at \url{https://huggingface.co/lamm-mit/x-lora}. The X-LoRA-Gemma model is available at \url{https://huggingface.co/lamm-mit/x-lora-gemma-7b}. The \texttt{Mistral.rs} inference engine is provided at \url{https://github.com/EricLBuehler/mistral.rs}.

\subsection{Mathematical formulation of X-LoRA} We define the $\mathbf{scalings}$ as a matrix comprising token-level and layer-level coefficients for LoRA adapters. The X-LoRA $\mathbf{scaling}$ $\mathbf{head}$ processes the hidden states from the base model via linear fully-connected layers to predict scalings, followed by a softmax operation along the last dimension to ensure that the scalings sum to one across all experts. ReLU activation functions and dropout layers are interleaved between the linear fully-connected layers within the X-LoRA scaling head. The hidden layer size is configurable and can be selected to effectively reduce the dimensionality from the hidden dimension to the scaling dimension or to implement a widening-and-narrowing scheme akin to the feed-forward layer in a transformer block.

The scalings form a matrix $\Lambda \in \mathbb{R}^{s \times l \times n}$, where $s$, $l$, and $n$ denote the sequence length, number of layers, and number of adapters, respectively. For each layer, the relevant scalings are extracted as a matrix $\lambda \in \mathbb{R}^{s \times n}$. These scalings $\lambda$ are then applied to each LoRA adapter via broadcast multiplication of the respective X-LoRA scaling $\lambda_i \in \mathbb{R}^{1 \times 1 \times s}$ to the inputs of the decomposition matrices $A$ and $B$.

With $W_0 \in \mathbb{R}^{d \times k}$, $B_i \in \mathbb{R}^{d \times r_i}$, and $A_i \in \mathbb{R}^{r_i \times d}$, where the rank $r_i \ll \min(d,k)$ is specific to each adapter $i$, the forward pass of an X-LoRA layer is given by:

\begin{equation}
    h = W_0x + \sum_{i=1}^{n} B_iA_i(x \cdot  \lambda_i) \alpha_i
\end{equation}

Visual representations of the scalings are shown in the main paper, depicted either through bar plots or via the \texttt{contourf} function in Matplotlib~\cite{MatplotlibDocumentation}.

\subsection{X-LoRA implementation algorithm and code} The implementation presented here prioritizes user-friendliness and the flexibility to readily extend the method to various models, particularly those within the Hugging Face framework. The X-LoRA algorithm functions by:

\begin{enumerate}
    \item Keeping the base model and adapters fixed while tuning the performance.
    \item Combining LoRA adapters based on the predictions made by the X-LoRA scaling head.
\end{enumerate}

All code is developed using Python and PyTorch~\cite{Paszke2019PyTorchLibrary}. 

\subsubsection{Dual forward pass strategy for self-aware inference} Since X-LoRA necessitates computing hidden states to estimate the scalings, we adopt a dual forward pass technique where the first pass allows the model to analyze the task and determine how to adjust itself, followed by the inference pass. This introduces a form of ‘self-awareness’, enabling X-LoRA to allocate extra computation for reflection instead of immediate output.

We refer to the initial pass as the scaling pass and the subsequent pass as the forward pass. During the scaling pass, the scalings $\Lambda$ are set to zero. We also tested assigning $\Lambda$ a value of ${n}^{-1}$ and observed comparable outcomes, suggesting stable performance (this parameter can be configured in the X-LoRA package). To facilitate applying X-LoRA across different models, we incorporate a forward pass hook. This mechanism is crucial for managing the interplay between the scaling and forward passes and constitutes the main feature that permits compatibility with all models.

\begin{enumerate}
    \item $\mathbf{Scaling}$ $\mathbf{pass}$: Executed by running the base model with adapters fixed at constant values. The hidden states produced are then used by the X-LoRA scaling head to generate $\Lambda$.
    \item $\mathbf{Forward}$ $\mathbf{pass}$: Utilizes $\Lambda$ to blend the adapters during the model’s forward pass. The resulting logits are outputted as the X-LoRA model’s final result.
\end{enumerate}

A cautionary note: employing the same key-value cache for both the scaling pass and forward pass can disrupt correct input to the scaling head because the cache persists and accumulates over the two passes. Therefore, we avoid using the key-value cache during training or inference and have implemented proper management of this behavior in \texttt{Mistral.rs}.

\subsubsection{Freezing weights and selectively enabling segments of the model for training} In an X-LoRA model, the base model and adapters remain fixed ({\it i.e.}, their parameters are not updated during training), making the X-LoRA scaling head the sole trainable component. The provided code also allows the option to unlock all adapters along with the X-LoRA scaling head for training. Furthermore, one could choose to jointly train the entire model — encompassing the X-LoRA scaling head, adapters, and the base foundational model — potentially employing distinct learning rates for each. Another potential direction for future research is to extend these ideas to design scaling functions bridging two or more foundational models, facilitating a traceable realization of SLERP-like modeling. This avenue remains open for exploration in subsequent work.

\subsubsection{X-LoRA layers} An X-LoRA layer is designed to scale and subsequently integrate the outputs from multiple adapters. It encapsulates the original LoRA layer, thereby granting access to the specific adapters assigned to that layer.

To enable the application of X-LoRA across arbitrary models, the algorithm replaces the original LoRA layers by constructing X-LoRA layers and embedding their forward computation logic into the LoRA layers. This approach utilizes features of Python’s object system to facilitate such layer replacement.

In essence, each X-LoRA layer performs the following operations during the forward pass (refer to Figure \ref{fig:general_arch}):

\begin{enumerate}
    \item $\mathbf{Scaling}$: Scale the input to each LoRA adapter by its corresponding scaling factor $\lambda_i$.
    \item $\mathbf{Forward}$: Execute the forward computation of the wrapped LoRA layer.
\end{enumerate}

\subsection{Optimized Inference with \texttt{Mistral.rs} implemented in Rust} To enhance inference throughput, we have developed \texttt{Mistral.rs}, a large language model serving framework written in Rust~\cite{matsakis2014rust}, which incorporates X-LoRA along with various performance-enhancing optimizations. The following techniques have been integrated, presented here in no specific order:

\begin{itemize}
    \item LoRA adapter weight stacking: When all LoRA adapter matrices $A_i$ and $B_i$ share identical dimensions across adapters, these weight matrices can be stacked along their first dimension.
\end{itemize}

Mathematically, given $W_0 \in \mathbb{R}^{d \times k}$, for adapter matrices $A_i \in \mathbb{R}^{r \times d}$ and $B_i \in \mathbb{R}^{d \times r}$ with $r_i \ll {\rm min} (d,k)$, we concatenate these matrices to form $A\in \mathbb{R}^{i \times r \times d}$ and $B\in \mathbb{R}^{i \times d \times r}$. The parameters $\alpha_i$ are also incorporated during the initialization of these stacked matrices.

To implement scalings, the dimensions are permuted such that $\lambda \in \mathbb{R}^{i \times 1 \times 1}$, and then applied via broadcast multiplication.

\item Non-granular scalings To minimize the number of times the scaling operation is executed, we provide an option for non-granular scalings. This method caches the scaling factors after $n$ completion cycles, leveraging the KV cache to maintain the sequence length at one for the rest of the request. This approach significantly lowers latency by reducing the frequency of base model invocations.

\item Fused Compute Unified Device Architecture (CUDA) kernels To enhance the efficiency of the forward pass, we utilize fused CUDA kernels. These fused kernels allow complex operations—such as Rotary Embedding (\cite{su2021roformer}), which would otherwise run sequentially—to be merged into a single kernel executable in parallel on a GPU. Our implementation draws from the vLLM framework (\cite{kwon2023efficient}).

\item Quantization The inference engine supports quantization, enabling the use of a quantized base model. Preliminary experiments show that quantization down to 8 bits performs effectively even with models initially trained in \texttt{bfloat16} format.

\end{itemize}

\subsection{Datasets}

Each adapter is trained on specialized datasets, which are summarized in Table~\ref{tab:table_loradata} along with citations to the original sources and literature.

\subsection{Training strategy} Training proceeds in multiple phases. Starting with the Zephyr-7B-$\beta$ model~\cite{HuggingFaceH4/zephyr-7b-betaFace}, which is built on top of the Mistral-7B model \cite{Jiang2023Mistral7B}, we train a series of adapters using the datasets detailed above. The Zephyr-7B-$\beta$ tokenizer is employed throughout.

The first eight adapters are trained on question-answer pairs directly from the datasets (see Table~\ref{tab:table_loradata}).

Protein mechanics tasks are trained in a similar fashion but involve specific forward and inverse instruction sets. The bidirectional instruction tasks include:

\begin{quote}
\tiny {\texttt{CalculateForce< G E E C D C G S P ....> [0.310]\\
CalculateEnergy< G E E C D C G S P ....> [0.220]\\ 
CalculateForceHistory< G E E C D C G S P ....> [0.004,0.034,0.125,0.142,0.159,0.102,0.079,0.073,0.131, ...]\\
GenerateForce<0.310>\,[ G E E C D C G S P ....]\\
GenerateEnergy<0.220>\,[ G E E C D C G S P ....]\\
GenerateForceHistory<0.004,0.034,0.125,0.142,0.159,0.102,0.079,0.073,0.131, ...>\,[ G E E C D C G S P ....] }}
\end{quote}

These correspond to three main tasks in total: determining the maximum unfolding force of a protein, calculating the unfolding energy, and obtaining the force-deformation history. Each of these forward tasks is paired with an inverse task of the same kind, which instead outputs a protein sequence, resulting in six tasks overall.

The rank assigned to each LoRA adapter is $r=16$ with a scaling factor of $\alpha=16$. The adapters are trained using five specific target modules: \texttt{v\_proj}, \texttt{k\_proj}, \texttt{q\_proj}, \texttt{gate\_proj}, and \texttt{o\_proj}, applying a dropout rate of 0.05. Typically, LoRA adapters are trained with low ranks ranging from 4 to 32 because it has been demonstrated that increasing the rank generally does not enhance performance but does raise computational costs. Lower ranks produce a more compact model with fewer parameters~\cite{hu2021lora}. The selected values for $r$, $\alpha$, and related settings follow commonly accepted standards in the literature that have proven effective. Furthermore, prior work by the author has explored various configurations and concluded that these values provide a reasonable balance~\cite{Luu2023BioinspiredLLM:Materials}.

The X-LoRA scaling head is trained using approximately 2,000 samples, which are a subset of the original dataset employed to train the adapters (we also augmented the training set with additional GPT-4 generated multiple choice question samples). We utilize a \texttt{paged\_adamw\_8bit} optimizer with a gradient clipping threshold of 0.3, a learning rate of $2\times 10^{-4}$ with warmup, and accumulate gradients over four steps, using a batch size of one. The X-LoRA model is trained for roughly 10,000 steps, achieving a loss near 0.28 by the end of training.

Given that the base Mistral architecture applied here consists of 32 layers, and adapters are incorporated into 5 transformer layers each, there are a total of 160 LoRA layers, each replaced by an X-LoRA layer. The X-LoRA scaling head predicts scaling factors for all 160 LoRA layers (for nine LoRA experts, this yields $160 \times 9 = 1440$ $\lambda_i$ values calculated for every token). The model employs a single linear layer with ReLU activations and dropout set at $p=0.2$, totaling around 5 million parameters. The code is designed with modularity to enable the creation and evaluation of more sophisticated deep classifier heads in other use cases.

Both LoRA adapters and the X-LoRA scaling head are trained through next-token prediction using cross-entropy loss (we estimate the likelihood that the model correctly predicts the subsequent token in a sequence based on the preceding tokens; the training objective is to maximize the probability of the correct next token as specified in the training dataset). We apply token shifting, where the input sequence is shifted by one token to generate the target sequence the model aims to predict. For instance, if the input sequence is ``Proteins are fundamental to", then the target sequence for training becomes ``are fundamental to biology". Consequently, the model attempts to predict the next token in the sequence.

As illustrated in Figure~\ref{fig:hierarchical_concept}, during the training of the X-LoRA scaling head, only the scaling head parameters are updated while all other model parameters remain fixed. The standard Zephyr tokenizer and prompt template are employed.

\subsection{X-LoRA-Gemma model} The X-LoRA-Gemma model is constructed similarly to the X-LoRA model described earlier, but it is based on the Gemma-7B-it model~\cite{Google/gemma-7b-itFace} and utilizes four adapters (refer to the main text for further details). These four adapters are trained on datasets covering mechanics and materials, protein mechanics, bioinspired materials, and an additional quantum-mechanics-based molecular properties dataset QM9 (see Table~\ref{tab:table_QM9})~\cite{Ruddigkeit2012EnumerationGDB-17,Ramakrishnan2015ElectronicSpace}.

Each LoRA adapter has a rank of $r=16$ and $\alpha=16$. Training is applied to the following seven target modules: \texttt{q\_proj}, \texttt{o\_proj}, \texttt{k\_proj}, \texttt{v\_proj}, \texttt{gate\_proj}, \texttt{up\_proj}, and \texttt{down\_proj}, with a dropout rate of 0.05. The rank for each adapter is fixed at $r=16$ with $\alpha=16$. Considering the four adapters and seven target modules, and the total of 28 layers in the Gemma-7B base model, a total of $28 \times 28 = 784$ $\lambda_i$ values are computed for each token.

The first two adapters (bioinspired and mechanics of materials) integrated into this model are trained using question-answer pairs as provided by their respective datasets (see Table~\ref{tab:table_loradata}).

Protein mechanics tasks are trained as described in the previous section, employing bidirectional forward and inverse instruction sets. Furthermore, an adapter trained on the QM9 dataset is included, with the following two tasks:

\begin{quote}
\tiny {\texttt{CalculateMolecularProperties<CC(C)N1CC1C=O> [0.098,0.358,0.581,0.395,0.309,0.330,0.570,0.649,0.519,0.519,0.519,0.518]\\
GenerateMolecularProperties<0.098,0.358,0.581,0.395,0.309,0.330,0.570,0.649,0.519,0.519,0.519,0.518> [CC(C)N1CC1C=O]\\ 
...
}}
\end{quote}

These encompass a total of two tasks: predicting molecular properties (the forward task) and generating a molecule with specified target properties (the inverse task). All 12 properties included in the QM9 dataset are either predicted or designed for, respectively (refer to Table~\ref{tab:table_QM9}).  
This model incorporates a total of four adapters.

The X-LoRA scaling head for this model consists of three layers, with the intermediate layer sized at $3072 \times {FF}_{\rm fac} = 12,288$, where ${FF}_{\rm fac}=4$ represents the feed-forward multiplier. The architecture employs linear layers paired with ReLU activation functions and a dropout rate of $p=0.2$, comprising approximately 47 million parameters.

The X-LoRA-Gemma scaling head was trained using roughly 17,000 samples. These were selected as a subset from the training dataset utilized for adapter training; it is important to highlight that although the X-LoRA-Gemma model contained only four adapters, the samples were drawn from the full training set of the aforementioned model, supplemented with samples from the QM9 dataset. This approach was taken to explore whether the X-LoRA model might acquire additional capabilities during the scaling head training phase. Given that the loss converged well to 0.58 by the end of training, alongside strong performance even on tasks the model was not explicitly trained for (e.g., chemistry, demonstrated by the molecule design example in the paper), we believe this generally produces favorable outcomes.

We employed a \texttt{paged\_adamw\_8bit} optimizer with gradient clipping set to 0.3, a learning rate of $1\times 10^{-4}$ incorporating warmup, and four steps of gradient accumulation, using a batch size of one. The X-LoRA-Gemma model was trained for approximately 62,000 iterations.

The native Gemma tokenizer along with its prompt template was utilized.

\subsection{Adversarial agentic modeling} Our adversarial agentic approach involves creating two X-LoRA agents. One agent specializes in formulating questions, initiating with an initial inquiry and continuously generating follow-up questions to uncover weaknesses and challenges within the responses. The other agent is responsible for answering these questions and providing ongoing replies. This methodology is implemented in \{Guidance\}~\cite{Guidance-ai/guidance:Models.}.

For the example that connects music and mechanics mentioned earlier, the characteristics of the two agents are defined as follows. Firstly, the physicist and questioner:

\begin{quote}
\small\texttt{You are a critical physicist participating in a discussion from the perspective of physics. Keep your responses concise and always challenge statements provocatively.

As a creative thinker, you incorporate ideas from other disciplines and push conceptual boundaries.}
\end{quote}

Secondly, the philosopher (in certain examples, this role is adapted to represent a biology expert):

\begin{quote}
\small\texttt{You are a philosopher knowledgeable in biology, chemistry, and mathematics. You are engaged in a discussion.

Keep your replies brief, yet accurate and imaginative. You generate excellent ideas and novel thought directions, always maintaining logical consistency.}
\end{quote}

The questioner is explicitly directed to reflect on the question and previous answers, and to develop new, insightful queries. This is achieved through a prompt structured as follows:

\begin{quote}
\small\texttt{Examine the following question and answer.

\#\#\# Question: <question>

\#\#\# Response: <response>

\#\#\# Instruction: Provide a SINGLE follow-up question that critically examines the response.

Do NOT provide an answer or commentary at this stage.

The single question is: }
\end{quote}

After the dialogue has unfolded for $N$ exchanges, the model’s task is to 1) summarize the main insights discussed, 2) enumerate the key insights as bullet points, and 3) pinpoint the most significant conclusion.

The unprocessed conversation text serves as the foundation for subsequent analysis through graph creation.

\subsection{Knowledge graph generation} We employ Zephyr-7B-$\beta$ to extract triplets from text, following the method outlined in~\cite{Rahulnyk/knowledge_graph:QnA} with enhancements inspired by the Llama Index graph construction procedure~\cite{Liu_LlamaIndex_2022}. The graphs are displayed using NetworX\cite{Networkx/networkx:Python} and Pyvis~\cite{WestHealth/pyvis:Graphs.}. A primary instruction, based on the strategy recommended in~\cite{Liu_LlamaIndex_2022}, is:

\begin{quote}
\small\texttt{Your objective is to derive the ontology of terms mentioned within the provided context. These terms should capture the essential concepts relevant to the context.

Present your output as a JSON list. Each list item includes a pair of terms and the relation between them, formatted like the following: [...] }
\end{quote}

We instruct the construction of triplets comprising nodes and edges as follows, consistent with the approach in~\cite{Liu_LlamaIndex_2022}:

\begin{quote}
\small\texttt{\begin{itemize}
            \item "node\_1": "A concept extracted from the ontology"
            \item "node\_2": "A related concept from the ontology"
            \item "edge": "A concise description of the relationship between node\_1 and node\_2"
        \end{itemize}
        }
\end{quote}

We also provide the model with examples of ontological analysis, such as:

\begin{quote}
\small\texttt{
       Context: ```Spider silk is a strong natural fiber used to catch prey in a web.``` }
\end{quote}

Followed by example triplets like

\begin{quote}
\small\texttt{\begin{itemize}
            \item "node\_1": "spider silk"
            \item "node\_2": "fiber"
            \item "edge": "is"
        \end{itemize}
        }
\end{quote}

The generated graph data are then aggregated and partitioned into communities utilizing the Girvan-Newman algorithm~\cite{Girvan2002CommunityNetworks}.

\subsection{Visualization of molecular structures}

We utilize PyMOL~\cite{Schrodinger2015The1.8} to visualize and examine the predicted protein structures. Various scripts and functions within this tool assist in identifying features such as secondary structures, hydrophobic and hydrophilic regions, disulfide bridges, and hydrogen bonds.

\section*{Data Availability Statement} The code and datasets underpinning the findings of this research are publicly accessible at \url{https://github.com/EricLBuehler/xlora} and \url{https://huggingface.co/lamm-mit/x-lora}. Additional data sources referenced in this study can be found, specifically in Table~\ref{tab:table_loradata}.

\end{document}